\subsection{Reptile vs. FOMAML: comparable performance, computational trade-offs}
\label{sec:disc:fomaml}

Reptile and FOMAML achieve statistically comparable $F_1$ scores across both random and spatial-CV protocols, with bootstrap 95\% confidence intervals that overlap in the majority of target-$K$ cells (Table~\ref{tab:spatial_results}). At $K=1$ under spatial CV, FOMAML wins in Copiap\'o ($0.683$ vs.\ Reptile $0.665$, $+1.8$\,pp), while Reptile wins in Huasco ($0.755$ vs.\ $0.745$, $+1.0$\,pp), Elqui ($0.476$ vs.\ $0.472$, $+0.4$\,pp), and Limar\'i ($0.528$ vs.\ $0.514$, $+1.4$\,pp); none of these inter-method differences exceed the bootstrap standard error. This convergence is consistent with \citet{raghu2020rapid}, who argue that the principal benefit of MAML-family methods is rapid feature reuse, not the second-order curvature term that distinguishes the two algorithms. Practically, Reptile's update rule (average-of-inner-displacements) is computationally cheaper than FOMAML's (requires a query-loss forward pass through the inner-adapted model per outer step), so we recommend Reptile as the default meta-learning algorithm for tabular landslide susceptibility transfer, reserving FOMAML for settings closer to Copiap\'o (hyper-arid, very small inventory) where the slight $K=1$ advantage may compound across applications.

This near-equivalence extends to a third meta-learner: ProtoNet \citep{snell2017protonet}, a metric-based method that performs \emph{no} gradient adaptation at test time, matches or exceeds the gradient-based methods at $K=1$ (Section~\ref{sec:results:extended}, Table~\ref{tab:extended_baselines}). The convergence of three structurally distinct meta-learners---averaged-displacement (Reptile), query-gradient (FOMAML), and nearest-prototype (ProtoNet)---onto similar performance strongly supports the \citet{raghu2020rapid} thesis that the few-shot advantage is driven by a transferable feature representation rather than by any particular adaptation rule. The failure of Meta-Baseline \citep{chen2021metabaseline}, which freezes a concat-pre-trained encoder, indicates that this representation must be learned \emph{episodically} (or refined by gradient adaptation): a representation optimized only for source classification does not transfer as well as one optimized for rapid task adaptation.

\subsection{Why DANN fails in the smallest target basin}
\label{sec:disc:dann}

DANN does not improve over Independent on average. The under-performance is concentrated in Copiap\'o, the smallest target basin in our experiment ($N=19$ landslide events), where DANN trails Independent by 5--7\,pp $F_1$ across all $K$ values (Table~\ref{tab:spatial_results}). In the other three target basins (Huasco, Elqui, Limar\'i), DANN is roughly neutral or marginally better than Independent (within $\pm 3$\,pp). This basin-dependent pattern contradicts the narrative in much of the DA literature, where DANN is presented as a strong default baseline, and demands explanation. We propose two complementary mechanisms.

First, \emph{adversarial invariance requires sufficient domain diversity}. With only $N=10$ source basins, the domain classifier easily exploits low-level features such as climate (bio\_12 ranges from $\sim 1$ mm/yr in the Atacama sources to $\sim 2{,}630$ mm/yr in Patagonian Magallanes) or lithology (volcanic vs.\ sedimentary). The Gradient Reversal Layer encourages the feature extractor to discard these signals, but they are also potentially predictive of landslide susceptibility (e.g., precipitation-driven debris flows correlate with bio\_12). The result is a representation that is more domain-invariant but less class-discriminative---a form of \emph{negative transfer} where alignment removes signal.

Second, \emph{the $K$-shot adaptation step is poorly matched to DANN's design philosophy}. DANN was designed for the unsupervised setting in which target labels are entirely unavailable. Re-purposing the trained DANN as a starting point for a small fine-tune on $K$ labels conflates two mechanisms (adversarial invariance + supervised fine-tune) that are not jointly optimized. Two alternative formulations---MMD-based DA \citep{long2015dan} and source-target joint embeddings---are reasonable next baselines but are outside our current scope.

\subsection{Spatial heterogeneity in Elqui: a hard transfer case}
\label{sec:disc:elqui}

Among target basins, Elqui exhibits the largest drop in $F_1$ between random and spatial CV protocols (27~pp at $K=10$; Table~\ref{tab:spatial_results} vs.\ Figure~\ref{fig:auc_random}). All methods, including FOMAML, are constrained to $F_1 \approx 0.50$ under spatial CV at any $K$. This indicates that intra-basin heterogeneity, not the cross-basin meta-learning capability, is the dominant factor in Elqui.

Two hypotheses are consistent with the data. (i) \textbf{Orographic structure}: Elqui has a strong west-east elevation gradient (sea level to $\sim 5{,}000$ m), with the cordillera dominated by Mesozoic plutonic rocks (granitoids), the precordillera by Cenozoic volcanics, and the coast by Quaternary alluvial deposits. Each lithological unit produces distinct landslide styles (debris flows in the alluvium, rotational slides in the granitoids), and our 11-feature representation may not capture these axes adequately. (ii) \textbf{Trigger heterogeneity}: the SERNAGEOMIN \emph{catastro} for Elqui aggregates events triggered by the 2015 Atacama floods, the 1997 El Ni\~no, and the 1922 Vallenar earthquake, three regimes with distinct preparatory factors. A meta-model trained on Maule 2010 (seismic) and Cha\~naral (rainfall) sources cannot easily recover this trigger-specific structure from the support set alone.

Future work should address Elqui-style heterogeneity by (a) richer spatial features (e.g., multi-temporal Sentinel-1 SAR coherence), (b) trigger-aware meta-learning, or (c) fold-specific stratification. None of these are straightforward extensions, and we report Elqui as a transparent failure mode rather than attempting to mask it.

\subsection{Climate robustness}
\label{sec:disc:climate}

The lack of correlation between source-target climate distance and transfer advantage (Section~\ref{sec:results:climate}) is a useful finding for practitioners. It suggests that meta-learning's prior captures generic geomorphic ``susceptibility patterns'' (e.g., concave hillslopes, high MRVBF, low landform variability) that transfer across climate regimes, rather than climate-specific signals. This in turn implies that source-pool diversity \emph{breadth} matters more than source-target \emph{similarity}: training on basins spanning multiple climate regimes is preferable to training on climate-matched sources only.

\subsection{Implications for remote-sensing-driven susceptibility mapping}
\label{sec:disc:rs}

Three findings carry methodological weight specifically for the remote-sensing-based geohazard community. \emph{First}, the meta-learned prior is recovered from an entirely open, globally-accessible feature stack---Sentinel-2 L2A bands, MERIT DEM derivatives, and WorldClim climatology---without proprietary inputs or in-situ measurements. This makes the framework directly replicable in any region with Sentinel-2 coverage (effectively, everywhere on Earth at 10--20~m resolution since 2017). The Chilean focus is illustrative, not restrictive. \emph{Second}, CFS-selected features (Section~\ref{sec:data:cfs}) reveal that DEM-derived terrain and GLCM-texture descriptors dominate the discriminative signal, while Sentinel-2 spectral bands and WorldClim bioclimatics are not retained. This has practical consequences: a 30~m DEM of moderate quality (MERIT, Copernicus GLO-30, or AW3D30) is the binding constraint on framework deployment, not optical-sensor revisit frequency. Countries without dense optical archives still have global DEM access. \emph{Third}, the meta-learning advantage is largest at $K=1$, which corresponds to one verified mass-movement footprint visible in high-resolution Sentinel-2 or Planet imagery. This aligns with operational reality: a remote-sensing analyst can typically validate a single landslide via visual interpretation in $<10$ minutes, whereas mapping the $K \geq 50$ events required by per-basin training takes weeks. Meta-learning thus reduces the human-in-the-loop cost from \emph{inventory construction} (weeks) to \emph{seed-event validation} (minutes), a regime where rapid operational deployment becomes feasible. A natural extension is to incorporate multi-temporal Sentinel-1 InSAR coherence and amplitude features \citep{poggi2026sentinellandslide, biggs2026volcanomonitor, elliott2026deformation} to capture surface-deformation precursors that static optical-DEM stacks cannot, an avenue actively explored for global InSAR landslide monitoring \citep{guo2026globalmonitor} and likely to amplify the meta-learning advantage by injecting time-varying signals into the support set.

\subsection{Limitations}
\label{sec:disc:limits}

We acknowledge five limitations.

\paragraph{Source domain count} With 10 source basins, the meta-learning regime is at the lower end of typical few-shot benchmarks (Omniglot uses 1{,}623 classes, miniImageNet uses 64 source classes). Adding more basins would likely increase the meta-learning advantage. Adding basins from neighboring countries (Peru, Argentina, Bolivia, Ecuador) is a natural next step.

\paragraph{MLP backbone} The main benchmark uses a 3-layer MLP on point-wise tabular features. This is intentionally conservative: it ensures interpretability, avoids large-scale GPU requirements, and maps cleanly to the practical workflow of basin managers. The U-Net-style CNN baseline of Section~\ref{sec:results:cnn} quantifies the cost of this choice on Huasco: the spatial CNN improves $F_1$ by 6--11~pp across methods and $K$ values, and the meta-learning lift is partially absorbed by the spatial representation. The MLP results should therefore be read as a conservative lower bound on achievable accuracy. A multi-target U-Net meta-learning evaluation, including richer patch sizes and attention-based architectures, is the natural follow-up.

\paragraph{Static climate} Our climate features (WorldClim 2.1 bioclimatics) represent a 1970--2000 baseline. Climate-change-driven shifts in precipitation regimes are a known driver of changing landslide hazard but are absent from our model. Coupling with CMIP6 projections is straightforward in principle.

\paragraph{Negatives generation} Eight of fourteen basins lack explicit negative samples, requiring us to generate them by random sampling within the basin polygon outside a 500~m buffer of positives. This introduces an implicit assumption---that all unlabeled pixels are negatives---which is statistically defensible (landslide footprints are sparse) but methodologically simplistic. Positive-Unlabeled (PU) learning is a more rigorous alternative.

\paragraph{Single-trigger framing} Our inventory mixes seismic and rainfall-triggered events without explicit modeling of trigger type. Trigger-aware sampling and feature design would likely improve discrimination, particularly in mixed-trigger basins like Elqui.

\subsection{Practical recommendations for new-basin susceptibility modeling}
\label{sec:disc:practical}

For a researcher or basin manager facing a \emph{newly-studied} watershed with fewer than 10 documented landslide events, our findings support the following protocol:

\begin{enumerate}
    \item \textbf{Identify a meta-training pool} of at least 4 source basins from a diverse range of hydro-climatic regimes. Climate matching is not required (Section~\ref{sec:results:climate}).
    \item \textbf{Prefer FOMAML over Reptile} for tabular feature representations. Use Reptile if computational budget is severely constrained.
    \item \textbf{Consider classical machine learning} (logistic regression, random forest) as a strong baseline; in spatial CV, classical ML matches meta-learning $F_1$ when $K \geq 5$. The unique advantage of meta-learning is \textbf{zero-step deployment}: the meta-trained model is immediately usable for a new basin without retraining.
    \item \textbf{Avoid DANN} when source-domain count is small ($N \leq 10$). A simple concat-source fine-tune or independent training is at least as effective and substantially simpler.
    \item \textbf{Validate with spatial cross-validation}, not random hold-out. Random CV inflates absolute $F_1$ by 5--25~pp depending on intra-basin spatial autocorrelation; the \emph{relative} method ranking is preserved, but the \emph{absolute} numbers reported to stakeholders should come from spatial CV.
    \item \textbf{Adapt for $\leq 30$ gradient steps}. Meta-learned initializations reach near-peak $F_1$ within 0--5 steps; further adaptation provides diminishing returns and risks over-fitting on the small support set.
    \item \textbf{Document hard cases}. Basins with strong intra-basin heterogeneity (multiple lithologies, multiple triggers, strong elevation gradient) may resist transfer and require trigger-stratified or lithology-stratified analysis as a follow-up.
\end{enumerate}
